% ============================================================
% 理论力学 · 交互讲义 —— M3《两体问题与系外行星》note.tex
% ============================================================
\documentclass[UTF8,a4paper,11pt]{ctexart}
\usepackage{amsmath,amssymb}
\usepackage{geometry}
\usepackage{booktabs}
\usepackage{hyperref}
\geometry{margin=2.5cm}
\hypersetup{colorlinks=true,linkcolor=blue,urlcolor=blue}

\title{\textbf{两体问题与系外行星}}
\author{理论力学 · 交互讲义}
\date{\today}

\begin{document}
\maketitle

\begin{abstract}
本讲义处理两体问题：从质心条件与约化质量出发，导出两体 Kepler 定律，
计算恒星绕质心的摆动速度，并用多普勒频移反推行星质量——这是系外行星
径向速度法的核心。以太阳--木星与太阳--地球两个系统为例，定量计算
摆动速度与光谱位移，解释为什么探测类地行星需要 $0.1\,\mathrm{m/s}$ 精度。
全部解析结果由配套脚本 \texttt{qa/numeric\_check.py}（仅依赖 numpy）逐项验证。
本模块覆盖两体专题（作业-4 第 5 题）。
\end{abstract}

\tableofcontents
\newpage

% ============================================================
\section{作业映射}
\label{sec:map}

\begin{center}
\small
\begin{tabular}{cp{7cm}p{3cm}c}
\toprule
来源 & 题面简述 & 核心知识点 & 对应小节 \\
\midrule
两体专题 & 太阳--木星/地球系统：质心距离、摆动速度、多普勒位移、类地行星探测 &
  质心、约化质量、两体 Kepler、RV、多普勒 & \ref{sec:derive} \\
\bottomrule
\end{tabular}
\end{center}

% ============================================================
\section{模型与假设}
\label{sec:model}

\begin{enumerate}
  \item \textbf{两体假设}：系统仅含两个质点，相互作用为牛顿万有引力
        $\boldsymbol F=-GM_1M_2/r^2\,\boldsymbol e_r$。
  \item \textbf{质点近似}：恒星与行星视为质点（忽略大小与自转）。
  \item \textbf{圆轨道近似}：本模块取圆轨道（$e\approx0$），
        椭圆修正在讨论节提及。
  \item \textbf{非相对论多普勒}：$\Delta\lambda/\lambda=v/c$，$v\ll c$。
  \item \textbf{edge-on}：轨道倾角 $i=90^\circ$（$\sin i=1$），
        视向速度最大。
  \item \textbf{单位制}：SI 单位；数值验证用题给数据。
\end{enumerate}

% ============================================================
\section{推导}
\label{sec:derive}

\subsection{质心与约化质量}
\label{sub:com}

两质量 $M_1$（恒星）、$M_2$（行星）相距 $a$。质心条件
\begin{equation}
  M_1 r_1 = M_2 r_2,\qquad r_1+r_2=a,
  \label{eq:com}
\end{equation}
解得
\begin{equation}
  r_1=\frac{M_2}{M_1+M_2}\,a,\qquad
  r_2=\frac{M_1}{M_1+M_2}\,a.
  \label{eq:r1r2}
\end{equation}
$r_1$ 为恒星到质心的距离，$r_2$ 为行星到质心的距离。

引入约化质量 $\mu=M_1M_2/(M_1+M_2)$，两体问题等价于质量 $\mu$ 的
质点在中心势 $-GM_1M_2/r$ 中运动：
\begin{equation}
  \mu\,\ddot{\boldsymbol r}=-\frac{GM_1M_2}{r^2}\,\boldsymbol e_r.
  \label{eq:reduced}
\end{equation}

\subsection{两体 Kepler 定律}
\label{sub:kepler2}

两体绕质心的角速度 $\omega$ 相同。由牛顿第二定律（以恒星为例）：
$M_1\omega^2 r_1 = GM_1M_2/a^2$。代入 $r_1=M_2 a/(M_1+M_2)$：
\begin{equation}
  \omega^2\frac{M_2 a}{M_1+M_2}=\frac{GM_2}{a^2}
  \quad\Longrightarrow\quad
  \boxed{\;\omega^2=\frac{G(M_1+M_2)}{a^3}\;}
  \label{eq:kepler2}
\end{equation}
即两体 Kepler 第三定律：$\omega$ 由总质量 $M_1+M_2$ 与间距 $a$ 决定。
轨道周期 $P=2\pi/\omega=2\pi\sqrt{a^3/(G(M_1+M_2))}$。

\subsection{两体能量与 virial}
\label{sub:energy2}

两体系统的总能量（圆轨道）：
\begin{equation}
  E = -\frac{GM_1M_2}{2a}.
\end{equation}
动能 $T=\frac12 M_1 v_1^2+\frac12 M_2 v_2^2
=\frac12 \mu v_{\mathrm{rel}}^2$，
势能 $V=-GM_1M_2/a$。virial 定理 $2\langle T\rangle=-\langle V\rangle$ 成立。

\subsection{椭圆轨道的两体修正}
\label{sub:ellipse2}

一般椭圆轨道：半长轴 $a$，离心率 $e$。两体 Kepler 第三定律仍为
$P^2=4\pi^2 a^3/(G(M_1+M_2))$，与单体形式相同，仅 $M\to M_1+M_2$。
恒星到质心的最大距离为 $r_{1,\max}=M_2 a(1+e)/(M_1+M_2)$。
径向速度半振幅为
$K=\frac{M_2}{M_1+M_2}\frac{2\pi a}{P\sqrt{1-e^2}}\sin i$；
正弦形式 $v_r(t)=K\sin(2\pi t/P+\phi)$ 仅适用于圆轨道（$e=0$），
一般椭圆轨道的 $v_r(t)$ 非正弦。
在固定 $M_1,M_2,a$ 与 $\sin i\neq0$ 下，$e\to1$ 时 $K\to\infty$
（形式上的牛顿极限）。

\subsection{恒星摆动速度}
\label{sub:rv}

恒星绕质心做圆周运动，摆动速度（切向）为
\begin{equation}
  v_1=\omega r_1 = \omega\,\frac{M_2}{M_1+M_2}\,a.
  \label{eq:v1}
\end{equation}
代入 \eqref{eq:kepler2}：
\begin{equation}
  v_1 = \frac{M_2}{M_1+M_2}\sqrt{\frac{G(M_1+M_2)}{a}}
  = \frac{M_2}{\sqrt{M_1+M_2}}\sqrt{\frac{G}{a}}.
  \label{eq:v1b}
\end{equation}
当 $M_2\ll M_1$ 时，$v_1\approx (M_2/M_1)\sqrt{GM_1/a}$。

\subsection{多普勒频移与径向速度}
\label{sub:doppler}

非相对论极限下，光源视向速度 $v$ 引起光谱波长偏移
\begin{equation}
  \frac{\Delta\lambda}{\lambda}=\frac{v}{c}.
  \label{eq:doppler}
\end{equation}
edge-on 轨道（$\sin i=1$）时视向速度振幅等于摆动速度 $v_1$。
测量 $\Delta\lambda$ 即可反推 $v_1$，进而由 \eqref{eq:v1b} 反推
行星质量 $M_2$（已知 $M_1$、$a$、$P$）。

\subsection{太阳--木星系统（作业题 (a)(b)(c)）}
\label{sub:jupiter}

给定：$M_\odot=1.989\times10^{30}\,\mathrm{kg}$，
$M_J=1.898\times10^{27}\,\mathrm{kg}$，
$a_J=7.78\times10^{11}\,\mathrm m$，
$P_J=11.86\,\mathrm{yr}=11.86\times3.156\times10^7\,\mathrm s$。

\paragraph{(a) 质心距离。}
由 \eqref{eq:r1r2}：
\begin{equation}
  r_1=\frac{M_J}{M_\odot+M_J}\,a_J
  =\frac{1.898\times10^{27}}{1.989\times10^{30}+1.898\times10^{27}}\times7.78\times10^{11}
  \approx 7.42\times10^8\,\mathrm m.
\end{equation}
太阳半径 $R_\odot=6.96\times10^8\,\mathrm m$，比值
\begin{equation}
  \frac{r_1}{R_\odot}\approx\frac{7.42\times10^8}{6.96\times10^8}\approx1.07.
\end{equation}
$r_1>R_\odot$：太阳质心在太阳表面之外，太阳绕质心的``摆动''幅度
超过太阳半径。

\paragraph{(b) 角速度与摆动速度。}
由 \eqref{eq:kepler2}：
\begin{equation}
  \omega=\sqrt{\frac{G(M_\odot+M_J)}{a_J^3}}
  =\frac{2\pi}{P_J}\approx1.68\times10^{-8}\,\mathrm{rad/s}.
\end{equation}
摆动速度
\begin{equation}
  v_1=\omega r_1\approx1.68\times10^{-8}\times7.42\times10^8
  \approx 12.5\,\mathrm{m/s}.
\end{equation}

\paragraph{(c) 多普勒位移。}
取 $\lambda=500\,\mathrm{nm}$，$v_1\approx12.5\,\mathrm{m/s}$：
\begin{equation}
  \Delta\lambda=\lambda\frac{v_1}{c}
  =500\,\mathrm{nm}\times\frac{12.5}{3.0\times10^8}
  \approx 2.08\times10^{-5}\,\mathrm{nm}.
\end{equation}

\subsection{太阳--地球系统（作业题 (d)(e)）}
\label{sub:earth}

给定：$M_E=5.97\times10^{24}\,\mathrm{kg}$，
$a_E=1.496\times10^{11}\,\mathrm m$，$P_E=1\,\mathrm{yr}$。

\paragraph{(d) 重算。}
\begin{equation}
  r_1=\frac{M_E}{M_\odot+M_E}\,a_E
  \approx\frac{5.97\times10^{24}}{1.989\times10^{30}}\times1.496\times10^{11}
  \approx 4.49\times10^5\,\mathrm m\;(449\,\mathrm{km}).
\end{equation}
\begin{equation}
  \omega=\frac{2\pi}{P_E}\approx1.99\times10^{-7}\,\mathrm{rad/s},\qquad
  v_1=\omega r_1\approx 0.089\,\mathrm{m/s}\;(8.9\,\mathrm{cm/s}).
\end{equation}
与走路速度 $1.4\,\mathrm{m/s}$ 相比，地球引起的太阳摆动慢约 16 倍；
与蜗牛速度 $1.3\,\mathrm{cm/s}$ 相比，快约 7 倍。
一句话：地球拖动太阳的摆动速度仅 $9\,\mathrm{cm/s}$，
比人走路慢一个数量级，但比蜗牛快。

\paragraph{(e) 为什么需要 $0.1\,\mathrm{m/s}$ 精度。}
\begin{enumerate}
  \item 信号本身仅 $\sim 0.09\,\mathrm{m/s}$（地球），仪器噪声必须优于信号，
        否则无法分辨。
  \item 太阳表面噪声（米粒组织、$p$ 模振荡等）量级 $\sim 1\,\mathrm{m/s}$，
        远大于地球信号，需要长时间平均与建模扣除。
\end{enumerate}

% ============================================================
\section{解析结果}
\label{sec:result}

\begin{enumerate}
  \item 质心距离：$r_1=M_2 a/(M_1+M_2)$（\eqref{eq:r1r2}）。
  \item 两体 Kepler：$\omega^2=G(M_1+M_2)/a^3$（\eqref{eq:kepler2}）。
  \item 摆动速度：$v_1=\omega r_1$（\eqref{eq:v1}）。
  \item 多普勒：$\Delta\lambda/\lambda=v/c$（\eqref{eq:doppler}）。
  \item 木星：$r_1\approx7.42\times10^8\,\mathrm m$，$r_1/R_\odot\approx1.07$，
        $v_1\approx12.5\,\mathrm{m/s}$，$\Delta\lambda\approx2.08\times10^{-5}\,\mathrm{nm}$。
  \item 地球：$r_1\approx4.49\times10^5\,\mathrm m$，$v_1\approx0.089\,\mathrm{m/s}$。
\end{enumerate}

% ============================================================
\section{数值验证}
\label{sec:numeric}

全部数值实验由配套脚本 \texttt{qa/numeric\_check.py} 完成（仅依赖 numpy）。

\begin{table}[htbp]
\centering
\small
\begin{tabular}{p{5cm}cc}
\toprule
物理量 & 解析值 & 数值值 \\
\midrule
$r_1$（太阳--木星） & $7.42\times10^8\,\mathrm m$ & $7.41\times10^8\,\mathrm m$ \\
$r_1/R_\odot$（木星） & $1.07$ & $1.06$ \\
$v_1$（木星） & $12.5\,\mathrm{m/s}$ & $12.4\,\mathrm{m/s}$ \\
$\Delta\lambda$（木星，500 nm） & $2.08\times10^{-5}\,\mathrm{nm}$ & $2.07\times10^{-5}\,\mathrm{nm}$ \\
$r_1$（太阳--地球） & $4.49\times10^5\,\mathrm m$ & $4.49\times10^5\,\mathrm m$ \\
$v_1$（地球） & $0.089\,\mathrm{m/s}$ & $0.089\,\mathrm{m/s}$ \\
\bottomrule
\end{tabular}
\caption{数值验证汇总。全部数值取自 \texttt{qa/numeric\_check.py} 的实际输出。}
\label{tab:numeric}
\end{table}

% ============================================================
\section{讨论}
\label{sec:discuss}

\paragraph{约化质量与等效单体问题。}
两体问题通过质心分离与约化质量 $\mu=M_1M_2/(M_1+M_2)$
等价为单体中心力问题。这一技巧在天体力学、分子物理中广泛使用。
相对运动的轨道形状与单体 Kepler 问题相同，仅时间尺度由总质量决定。

\paragraph{质心在太阳表面之外。}
木星引起的太阳摆动幅度 $r_1\approx7.4\times10^8\,\mathrm m>R_\odot$，
太阳绕质心的轨道半径超过太阳半径，恒星的`` Reflex motion'' 可观。

\paragraph{类地行星的挑战。}
地球信号 $\sim0.09\,\mathrm{m/s}$，比太阳表面噪声 $\sim1\,\mathrm{m/s}$
小一个数量级，需要长时间积分与多普勒测量精度 $\sim0.1\,\mathrm{m/s}$。

\paragraph{系外行星探测方法。}
径向速度法（RV）是最成功的系外行星探测方法之一，发现了数千颗行星。
其精度要求随行星质量与轨道距离变化：热木星（$M_p\sim M_J$，$a\sim0.05\,\mathrm{AU}$）
信号 $\sim 50\,\mathrm{m/s}$，地球类似行星信号 $\sim 0.1\,\mathrm{m/s}$。
其他方法包括凌星法（transit）、直接成像、微引力透镜等。

\paragraph{径向速度曲线的形状。}
对圆轨道，$v_r(t)$ 为正弦函数，振幅 $K=v_1$（\eqref{eq:v1}）。对椭圆轨道，
$v_r(t)$ 非正弦，振幅为
$K=\frac{M_2}{M_1+M_2}\frac{2\pi a}{P\sqrt{1-e^2}}\sin i$。
实际观测中拟合 Kepler 轨道得到 $K$、$P$、$e$，再反推 $M_2\sin i$（最小质量），
不直接使用圆轨道公式 \eqref{eq:v1b}。
\paragraph{与作业题的对照。}
表~\ref{tab:numeric} 覆盖了作业题 (a)--(e) 的全部数值结果。

% ============================================================
\begin{thebibliography}{9}

\bibitem{goldstein}
H.~Goldstein, C.~P.~Poole, J.~L.~Safko,
\emph{Classical Mechanics}, 3rd ed., Pearson, 2002.
（两体问题、质心与约化质量。）

\bibitem{mayor}
M.~Mayor, D.~Queloz,
\emph{A Jupiter-mass companion to a solar-type star}, Nature, 1995.
（系外行星径向速度法的里程碑。）

\bibitem{zhou}
周衍柏，《理论力学教程》（第三版），高等教育出版社，2009.
（两体问题与有心力场。）

\end{thebibliography}

\end{document}
